\documentclass[11pt]{article}

% Basic packages
\usepackage[a4paper,margin=1in]{geometry}
\usepackage{amsmath,amssymb,amsthm,mathtools}
\usepackage{graphicx}
\usepackage{subcaption}
\usepackage{hyperref}
\usepackage{bookmark}
\usepackage{titlesec}
\usepackage{indentfirst}
\usepackage{natbib}
\usepackage{booktabs}
\usepackage{xcolor}
\usepackage{authblk}

% Language / encoding
\usepackage[utf8]{inputenc}
\usepackage[LGR,T1]{fontenc}
\usepackage[greek,english]{babel}

% Quotation handling
\usepackage{csquotes}

% Custom macros
%% ---------------------------------------------------------------
%%  macros.tex
%% ---------------------------------------------------------------

%% Editorial markers (remove before submission)
\newcommand{\TODO}[1]{\textcolor{red!70!black}{\textbf{[TODO: #1]}}}
\newcommand{\NOTE}[1]{\textcolor{blue!60!black}{\textit{[Note: #1]}}}

%% Dirac notation
\newcommand{\ket}[1]{\left|#1\right\rangle}
\newcommand{\bra}[1]{\left\langle#1\right|}
\newcommand{\braket}[2]{\left\langle#1\middle|#2\right\rangle}

%% Greek Aristotelian terms (requires babel greek + LGR encoding)
\newcommand{\dunamis}{\textgreek{δύναμις}}
\newcommand{\energeia}{\textgreek{ἐνέργεια}}
\newcommand{\entelechia}{\textgreek{ἐντελέχεια}}

\title{\textbf{Ontological Collapse: What the Classical Atom Reveals About Scientific Models and the Metaphysics They Carry}}
% \title{\textbf{Ontological Collapse: On Electrons, Atoms, Observers, and the Metaphysics they Carry}}
% \title{\textbf{Ontological Collapse: On the Hidden Metaphysics of Physical Models}}

\author[1]{Jes\'{u}s Urtasun Elizari}
\author[2]{Alberto Miguel G\'{o}mez}

\affil[1]{\small MRC LMS, Imperial College London,
          \texttt{jurtasun@imperial.ac.uk}}
\affil[2]{\small Departamento, Universidad,
          \texttt{email@institucion.ac.uk}}

\date{\today \\ \small Draft for discussion --- do not circulate \\ De paseos, perogrulladas y discusiones varias}

\begin{document}

\maketitle

\begin{abstract}
We present a discussion on the classical calculation of the radiative collapse of the hydrogen atom, and how it constitutes not merely an empirical failure but an \emph{ontological} one. 
The planetary model of the atom fails not because its equations are internally inconsistent, but because the type of entities it presupposes---such as a point charge with a definite trajectory---does not correspond to what electrons are. 
We use this case to distinguish two layers present in any physical model, often ignored or not sufficiently emphasized: its predictive mathematical content and its implicit ontological commitments. 
We then argue that Aristotle's distinction between potentiality (\textgreek{δύναμις}) and actuality (\textgreek{ἐνέργεια}), recovered by Heisenberg in 1958, provides a more adequate ontological framework for the quantum atom than either naive particle realism or strict instrumentalism. 
The paper closes by situating this reading within contemporary debates on structural realism and the metaphysics of quantum objects.
\end{abstract}

% At the very start of section1.tex, before subsection 1.1,
% replace the current first sentence with:

This paper argues that physical models can fail in two distinct ways:
predictively, when their equations give wrong numerical answers within
their domain, and ontologically, when the kind of entity they
presuppose does not exist at the relevant scale.
The classical calculation of the radiative collapse of the hydrogen
atom is our primary case study---a derivation that is mathematically
exact, internally consistent, and empirically false---from which we
derive a general diagnostic applicable to any physical model.
The same structure appears in the ultraviolet catastrophe and in the
quantum measurement problem, as we note in Section~2.
We then argue that Aristotle's distinction between potentiality and
actuality, recovered by Heisenberg in 1958, provides a more adequate
ontological framework for the quantum atom than either classical
particle realism or strict instrumentalism.

% Section .............................................................................................................
\section{The Classical Collapse: An Explicit Calculation}
\label{section:collapse}

The failure of the planetary models of the atom is a well known and extensively discussed example. We present the calculation in full because the philosophical argument that follows depends on it: the failure is not a rough estimate but an exact consequence of three internally consistent and independently verified physical laws.

% Subsection ..........................................................................................................
\subsection{The Three Ingredients and Their Sources}
\label{subsection:ingredients}

\textbf{Centripetal acceleration (Newton, 1687)}

\medskip

Newton's \emph{Principia Mathematica} \cite{newton_1687_principia} establishes that for uniform circular motion of radius $r$ and speed $v$, differentiating the position vector $\vec{r}(t) = r(\cos\omega t,\,\sin\omega t)$ twice with $\omega = v/r$ gives
\begin{equation}
  \ddot{\vec{r}} = -\omega^2\,\vec{r} \quad\Longrightarrow\quad \boxed{a = \frac{v^2}{r}} \; ,
  \label{eq:centripetal}
\end{equation}
directed toward the centre. No approximation is involved.

\medskip

\textbf{Coulomb's law (Coulomb, 1785)}

\medskip

Measured by torsion balance in 1785 \cite{coulomb_1785_electricity}, the force arising between two point charges $q_1$ and $q_2$ separated by distance $r$ is
\begin{equation}
  \boxed{F = \frac{1}{4\pi\epsilon_0}\frac{q_1\,q_2}{r^2}} \; ,
  \label{eq:coulomb}
\end{equation}
attractive for opposite charges. The constant $\epsilon_0$ would be added by Maxwell's 1865; the inverse-square dependence is Coulomb's direct experimental result.

\medskip

\textbf{Larmor radiation formula (Larmor, 1897)}

\medskip

From Maxwell's equations, the electric field radiated by an accelerated charge falls as $1/r$, so the Poynting flux integrated over a distant sphere yields a finite radiated power \cite{larmor_1897_radiation}:
\begin{equation}
  \boxed{P = \frac{e^2\,a^2}{6\pi\epsilon_0\,c^3}} \; .
  \label{eq:larmor}
\end{equation}
The dependence on $a^2$ follows because the radiated field is linear in $a$, and power is quadratic in the field amplitude.

% Subsection ..........................................................................................................
\subsection{The Differential Equation for \texorpdfstring{$r(t)$}{r(t)}}
\label{subsection:ode}

Setting Coulomb's force equal to the centripetal requirement for the electron (mass $m_e$, charge $e$) one obtains:
\begin{equation}
  \frac{e^2}{4\pi\epsilon_0\,r^2} = \frac{m_e\,v^2}{r} \implies v^2 = \frac{e^2}{4\pi\epsilon_0\,m_e\,r} \; .
  \label{eq:virial-v}
\end{equation}

The total mechanical energy is then:
\begin{equation}
  E = \frac{1}{2}m_e v^2 - \frac{e^2}{4\pi\epsilon_0\,r} = -\frac{e^2}{8\pi\epsilon_0\,r} \; ,
  \label{eq:energy}
\end{equation}
where the virial theorem ensures $E_{\mathrm{kin}} = -E/2$.

\medskip

As the electron loses energy by radiation, $E$ decreases and $r$ decreases monotonically. Substituting $a = e^2/(4\pi\epsilon_0 m_e r^2)$ into~(\ref{eq:larmor}):
\begin{equation}
  P = \frac{e^6}{12\pi^3\epsilon_0^3\,m_e^2\,c^3\,r^4}.
  \label{eq:power-r}
\end{equation}

Equating total energy $\dot{E} = -P$ and differentiating~(\ref{eq:energy}) one obtains:
\begin{equation}
  \frac{e^2}{8\pi\epsilon_0\,r^2}\,\dot{r} = -\frac{e^6}{12\pi^3\epsilon_0^3\,m_e^2\,c^3\,r^4} \; .
\end{equation}

Rearranging:
\begin{equation}
  r^2\,\mathrm{d}r = -\frac{2e^4}{3\,(4\pi\epsilon_0)^2\,m_e^2\,c^3}\,\mathrm{d}t \;\equiv\; -C\,\mathrm{d}t \; .
  \label{eq:ode}
\end{equation}

% Subsection ...........................................................................................................
\subsection{Integration and Result}
\label{subsection:result}

Integrating~(\ref{eq:ode}) from $r_0 = a_0$ (the Bohr radius, $a_0 \approx 0.529\,\text{\AA}$) to $r = 0$:
\begin{equation}
  \int_{a_0}^{0} r^2\,\mathrm{d}r = -C\,t_{\mathrm{collapse}} \implies \frac{a_0^3}{3} = C\,t_{\mathrm{collapse}},
\end{equation}

giving the final results
\begin{equation}
  \boxed{t_{\mathrm{collapse}} = \frac{(4\pi\epsilon_0)^2\,m_e^2\,c^3\,a_0^3}{4\,e^4} \approx 1.6\times10^{-11}\,\mathrm{s} \; .}
  \label{eq:tcollapse}
\end{equation}

Sixteen picoseconds. This is not an order-of-magnitude estimate; it is the exact result of integrating three independently verified and internally consistent physical laws. The calculation is impeccable, yet the conclusion is empirically false. That tension is the starting point of everything that follows.

\begin{table}[h]
  \centering
  \begin{tabular}{|l|l|l|}
    \toprule
    Constant & Symbol & Value \\
    \midrule
    Bohr radius            & $a_0$        & $0.529\times10^{-10}$\,m \\
    Electron mass          & $m_e$        & $9.109\times10^{-31}$\,kg \\
    Elementary charge      & $e$          & $1.602\times10^{-19}$\,C \\
    Speed of light         & $c$          & $2.998\times10^{8}$\,m\,s$^{-1}$ \\
    Permittivity of vacuum & $\epsilon_0$ & $8.854\times10^{-12}$\,F\,m$^{-1}$ \\
    \bottomrule
  \end{tabular}
  \caption{Physical constants used in the collapse calculation.}
  \label{tab:constants}
\end{table}

% Section .............................................................................................................
\section{Two Layers in Every Physical Model}
\label{section:layers}

The collapse calculation exemplifies a general structure present in all physical models, but rarely made explicit. Every model carries at least two distinguishable layers of content:

\begin{enumerate}
  \item \textbf{A predictive-mathematical layer:} the equations, boundary conditions, and inference rules that generate quantitative predictions about measurable quantities.
  \item \textbf{An ontological layer:} the implicit commitments about what kinds of entities exist, what properties they possess independently of measurement, and how they relate causally.
\end{enumerate}

Standard empiricist philosophy of science focuses almost exclusively on the first layer \cite{vanfraassen_1980_image}. A model is adequate if its predictions agree with experiment; inadequate otherwise. On this view, the classical atomic model fails because it predicts a collapse timescale of $\sim10^{-11}$\,s while atoms are stable indefinitely. That diagnosis is correct, but incomplete.

% Subsection ..........................................................................................................
\subsection{The Ontological Layer of the Classical Model}
\label{subsection:layers}

The deeper diagnosis is that the classical model fails at the second layer. Newton's equations for dynamics, Coulomb's law for electromagnetism, and Larmor's radiation formula are not wrong in themselves---they remain excellent approximations for macroscopic charged objects. Yet, they fail for the electron in the hydrogen atom because they embed a specific ontology:

\begin{itemize}
  \item Electrons are \emph{point particles} with definite position and
    momentum at every instant.
  \item They move along \emph{continuous trajectories} governed by
    deterministic force laws.
  \item They act as \emph{classical sources} for the electromagnetic
    field, coupling through their charge and acceleration.
\end{itemize}

It is known nowadays that each of these assumptions is false. A model can therefore be \emph{mathematically consistent}, \emph{empirically adequate in its domain}, and \emph{ontologically erroneous}. The classical hydrogen atom is a case study in exactly this combination.



% Subsection ..........................................................................................................
\subsection{The General Diagnostic Distinction}
\label{subsection:distinction}

We propose the following general principle:

\begin{quote}
  \emph{When a physical model fails empirically at the boundary of its domain, one must ask whether the failure is predictive or ontological. A predictive failure can often be corrected by modifying equations within the same ontological framework. An ontological failure requires replacing the picture of what exists.}
\end{quote}

As an example, the transition from Newtonian to relativistic mechanics is largely \emph{predictive}: the ontology of point particles with definite trajectories is preserved; only force laws and the metric of spacetime change. The transition from classical to quantum mechanics is \emph{ontological}: the very category of ``particle with a definite trajectory'' ceases to apply.

\medskip

This distinction has implications for how we interpret theoretical progress. Kuhnian paradigm shifts \cite{kuhn_1962_structure} and Lakatosian research programmes \cite{lakatos_1978_methodology} both capture aspects of it, but neither systematically separates the predictive from the ontological layer. The collapse calculation makes the separation vivid: the predictive machinery is flawless; it is the ontological presuppositions that must go.

\medskip

\TODO{Expand: Cartwright on laws vs.\ capacities; Giere on model-world relations; Chakravartty on semi-realism. Key point: none of these frameworks isolates the ontological layer as a separate diagnostic target in quite the way we propose.}

% Section .............................................................................................................
\section{Bohr's Provisional Resolution and the Persistence of the Problem}
\label{section:bohr}

% Subsection ..........................................................................................................
\subsection{The 1913 Postulates as Ontological Stipulation}

Bohr's response in 1913 was philosophically unusual in its candour \cite{bohr_1913_constitution}. He did not attempt to repair the classical ontology. He declared classical electrodynamics inapplicable to intra-atomic orbits and introduced quantised stationary states as irreducible postulates: the electron in a stationary state does not radiate, not because some mechanism suppresses radiation, but because the very concept of radiating simply does not apply to that state.

\medskip

This is an ontological statement, not a dynamical explanation. Bohr acknowledged as much, and indeed he described his postulates as provisional, awaiting a deeper theory. The immediate success was spectacular. The energy levels
\begin{equation}
  E_n = -\frac{m_e e^4}{2\hbar^2}\cdot\frac{1}{n^2} = -\frac{13.6\,\mathrm{eV}}{n^2}, \qquad n = 1, 2, 3, \ldots
  \label{eq:bohr-levels}
\end{equation}

reproduce the Balmer series \cite{balmer_1885_spectrallinien} and the full Rydberg formula \cite{rydberg_1888_linespectra}, with the Rydberg constant now derived rather than fitted:
\begin{equation}
  R_H = \frac{m_e e^4}{4\pi\hbar^3 c} \approx 1.097\times10^7\,\mathrm{m}^{-1}.
\end{equation}

The precision was stunning. Yet, the ontological justification was absent.

% Subsection ..........................................................................................................
\subsection{Quantum Mechanics: Relocation, Not Dissolution}

Full quantum mechanics relocates rather than dissolves the problem. In the Schr\"{o}dinger picture, the ground-state electron occupies the state $\psi_{1s}(r) \propto e^{-r/a_0}$, spherically symmetric. There is no definite orbit, no classical centripetal acceleration, and therefore no Larmor radiation. Stability is explained: the ground state is the lowest eigenvalue of the Hamiltonian $\hat{H} = \hat{p}^2/2m_e - e^2/4\pi\epsilon_0 r$, and there is no lower state to decay into.

\medskip

Quantum electrodynamics (QED) goes further, correctly computing excited-state lifetimes, the Lamb shift ($\sim$1057\,MHz for the $2s$--$2p_{1/2}$ splitting, measured by Lamb and Retherford in 1947 \cite{lamb_1947_finestructure}), and spontaneous emission rates at parts-per-billion precision.

% Subsection ..........................................................................................................
\subsection{The Tail That Remains Open}

Yet the problem has a residue. The self-energy of the electron in its own electromagnetic field has no fully satisfactory formulation. The Abraham--Lorentz equation \cite{abraham_1903_dynamik, lorentz_1904_electromagnetic, dirac_1938_radiating},
\begin{equation}
  m_e \dot{v} = F_{\mathrm{ext}} + \frac{e^2}{6\pi\epsilon_0 c^3}\ddot{v},
  \label{eq:AL}
\end{equation}
admits \emph{runaway solutions} (acceleration growing exponentially without external force) and \emph{pre-acceleration solutions} (particle responds before the force arrives, violating causality). In QED these pathologies are controlled by renormalisation, which works numerically with extraordinary precision but which Dirac regarded as a fundamental logical inconsistency \cite{dirac_1938_radiating} and Feynman called sweeping the problem under the rug with an elegant technique \cite{feynman_1985_qed}.

The ontological question Bohr implicitly raised in 1913---what \emph{kind of thing} is an electron, considered as both a quantum object and a source of its own field?---remains incompletely answered. Each theoretical layer has provided a more adequate partial ontology. None has provided a complete one.

\medskip

Consider a paragraph on effective field theory (Weinberg, Wilson): EFT explicitly embraces ontological incompleteness at each scale. That may be a more honest framing than ``the problem is unsolved.''

% Section .............................................................................................................
\section{Aristotle's Potentiality and the Quantum State}
\label{section:aristotle}

% Subsection ..........................................................................................................
\subsection{Heisenberg's Invocation}

In \emph{Physics and Philosophy} (1958), Heisenberg introduced an Aristotelian category to interpret quantum mechanics \cite{heisenberg_1958_physics}:

\begin{quote}
  ``[The probability wave] was a quantitative version of the old concept of \textit{potentia} in Aristotelian philosophy. It introduced something standing in the middle between the idea of an event and the actual event, a strange kind of physical reality just in the middle between possibility and reality.''
\end{quote}

This invocation is often treated as a philosophical metaphor appended to an otherwise self-contained physical theory. We argue it is more: it is a serious ontological proposal that deserves reconstruction and critical evaluation on its own terms.

% Subsection ..........................................................................................................
\subsection{The Aristotelian Framework}

Aristotle's \emph{Metaphysics} \cite{aristotle_1924_metaphysics} distinguishes two fundamental modes of being:

\begin{itemize}
  \item \dunamis\ (potentiality): what a thing can be or do, considered as a real but not yet actualised property. Not mere logical possibility, nor ignorance about an already determinate state, but a genuine ontological category---a real property of things, prior to and distinct from their actual states.
  \item \energeia\ / \entelechia\ (actuality): the realisation of that potentiality; what a thing is in full exercise of its nature.
\end{itemize}

For Aristotle, bronze \emph{is} potentially a statue before any sculptor acts; the seed \emph{is} potentially a tree. 
These are not descriptions of what we happen not to know; they describe what the thing really is, in a mode of being distinct from full actuality.

% -----------------------------------------------------------
% INSERT between the Aristotelian framework paragraph and
% "Application to the Quantum State"
% New subsection on Zubiri
% -----------------------------------------------------------

\subsection{A Note on Zubiri and the Etymology of Being}

The recovery of Aristotelian potentiality for physics is not without
precedent in the philosophical tradition, though it has gone largely
unnoticed in the analytic philosophy of physics.
Xavier Zubiri, in \emph{Sobre la esencia} \cite{zubiri_1962_esencia}
and \emph{Inteligencia sentiente} \cite{zubiri_1980_inteligencia},
developed a sustained critique of the ontological tradition precisely
through attention to the etymology and original sense of the Greek
\textgreek{ὄν} (\emph{on}, the present participle of \emph{einai},
to be).

Zubiri's central contention is that the Latin translation of
\textgreek{ὄν} as \emph{ens}---and the subsequent scholastic
hypostasization of \emph{esse} as an abstract noun---transformed
what was in Aristotle a \emph{dynamic} concept (being as an ongoing
process, a \emph{being-in-act}) into a \emph{static} one (being as
a fixed substantial nature).
The classical ontology of the atom inherits precisely this Latin
distortion: electrons are conceived as \emph{entia} with fully
determinate properties at every instant---a picture of being as
complete, static actuality.

What Zubiri calls \emph{la realidad} is not a fixed substrate of
properties but a \emph{de suyo}---a being-from-itself that is
irreducibly open, whose determinations emerge in relation to other
realities and cannot be read off from an intrinsic nature considered
in isolation.
This is structurally congruent with the potentiality reading: the
electron's \dunamis\ is not a deficient mode of being awaiting
completion, but a genuine and irreducible mode of \emph{de suyo}
that precedes and grounds any particular actualisation.

Zubiri's framework has not been applied to the philosophy of quantum
mechanics, to our knowledge.
We suggest it deserves to be: his critique of the static Latin
\emph{esse} maps precisely onto the failure of the classical atomic
ontology, and his recovery of dynamic \textgreek{ὄν} maps onto the
potentiality reading we are defending.
The convergence is not merely terminological---it reflects a shared
diagnosis that the metaphysical presuppositions inherited from
scholastic substance ontology are inadequate to describe entities
whose mode of being is irreducibly open and relational.

% Subsection ..........................................................................................................
\subsection{Application to the Quantum State}

The mapping is structurally precise. Before a position measurement on an electron in state $\psi$, the electron does not have a definite position we merely fail to know---this is what Bell's theorem and its experimental confirmation rule out \cite{bell_1964_epr, aspect_1982_bell, hensen_2015_loophole}.

\medskip

Rather, the electron has a distribution of \emph{potential} positions, encoded by the so-called wave function $|\psi(x)|^2$, which are real properties of the electron's present state. Measurement does not reveal a pre-existing value; it \emph{actualises} one of the potentialities.

\medskip

On this reading:
\begin{itemize}
  \item $\psi$ is the mathematical representation of the electron's potentialities.
  \item Measurement is the physical process of passage from potentiality to actuality.
  \item The Born rule $|\psi(x)|^2\,\mathrm{d}x$ gives the quantitative structure of the potentiality distribution.
\end{itemize}

This is not the Copenhagen interpretation in its instrumentalist form, which treats $\psi$ as a bookkeeping device for predictions. It is a realist reading: the potentialities are real features of the world, not features of our knowledge.

% Subsection ..........................................................................................................
\subsection{What This Resolves and What It Does Not}

The Aristotelian reading dissolves the puzzle of the classical collapse. The $1s$ electron has no definite trajectory; there is no classical kinematic state to insert into Larmor's formula~(\ref{eq:larmor}). The question ``why does the orbiting electron not radiate?'' does not arise, because its presupposition---that the electron has a definite centripetal acceleration---is false.

What this reading does \emph{not} resolve is the measurement problem in its general form. It relocates the hard question to the actualisation process: what determines which potentiality becomes actual? The major interpretations map partially onto different answers:

\begin{itemize}
  \item \textbf{Copenhagen:} actualisation is a primitive, theory-external fact; quantum mechanics describes potentialities but not the process of actualisation itself.
  \item \textbf{GRW:} spontaneous collapse provides a physical mechanism for actualisation, at the cost of modifying the Schr\"{o}dinger equation \cite{ghirardi_1986_unified}.
  \item \textbf{Everett:} all potentialities are actualised in branching worlds; no collapse occurs \cite{everett_1957_relativestate}.
  \item \textbf{Bohmian mechanics:} particles always have definite positions (full actuality); the pilot wave guides them deterministically \cite{bohm_1952_hiddenvariables}.
\end{itemize}

% Añadir tras el párrafo "We do not adjudicate between these here..."

\phantomsection\label{sec:wigner}

The measurement problem becomes sharper still in Wigner's
thought experiment of 1961 \cite{wigner_1961_friend}.
An observer inside an isolated laboratory---Wigner's
friend---measures the spin of an electron and, from his
perspective, obtains a definite outcome.
Wigner, outside the laboratory, describes the entire system
(friend, apparatus, electron) as a closed quantum system
evolving unitarily, assigning it the entangled state
\begin{equation}
  |\Psi\rangle = \frac{1}{\sqrt{2}}
  \bigl(|\uparrow\rangle_e|\mathrm{saw}\,\uparrow\rangle_F
      + |\downarrow\rangle_e|\mathrm{saw}\,\downarrow\rangle_F\bigr).
  \label{eq:wigner}
\end{equation}
The two observers assign incompatible quantum states to the same
physical system at the same moment.
On the potentiality reading, the paradox sharpens a question that
our framework does not fully answer: is the passage from
potentiality to actuality---from $\delta\acute{\upsilon}\nu\alpha\mu\iota\varsigma$
to $\epsilon\nu\acute{\epsilon}\rho\gamma\epsilon\iota\alpha$---an
objective physical event, or is it indexed to a particular
observer's interaction with the system?
If the former, when exactly does it occur?
If the latter, the potentialities are not fully ontic in the
sense we have been defending.

Frauchiger and Renner \cite{frauchiger_2018_quantum} sharpened
this into a no-go theorem: assuming simultaneously that quantum
mechanics is universal, that different observers' reasoning is
mutually consistent, and that the Born rule applies,
leads to contradiction.
The theorem does not refute any single interpretation but shows
that the three assumptions cannot all be true together---each
major interpretation must abandon at least one.
For the potentiality realist, the theorem suggests that the
ontology of actualisation requires further precision: a
fully satisfactory account must specify the physical conditions
under which joint potentialities become actual, and for whom.

\medskip

We do not adjudicate between these here. Our claim is more modest: the Aristotelian framework provides a more adequate \emph{starting} ontology for quantum mechanics than either classical particle realism or strict instrumentalism, and the collapse calculation is the sharpest way to see why.

\TODO{Reviewer will ask: is potentiality realism distinguishable from an epistemic interpretation? Key argument: potentialities are ontic, not epistemic---they ground the Born-rule probabilities rather than describing ignorance. See \cite{margenau_1954_interpretations} and \cite{bunge_1979_ontology}.}

% -----------------------------------------------------------
% INSERT replaces the \TODO at the end of section 4.4
% -----------------------------------------------------------

A natural objection must be addressed directly: is the potentiality
reading genuinely \emph{ontic}, or is it merely a reformulation of an
epistemic interpretation in Aristotelian vocabulary?
The objection runs as follows---the wavefunction $\psi$ describes our
knowledge of the system; ``potentiality'' is just a philosophically
dressed version of ``probability given our ignorance''; and the
Aristotelian language adds no ontological content.

We resist this reading on two grounds.

First, Bell's theorem \cite{bell_1964_epr} and its experimental
confirmation \cite{aspect_1982_bell,hensen_2015_loophole} rule out
any interpretation on which the wavefunction merely encodes ignorance
about pre-existing local values.
If $|\psi(x)|^2$ described our ignorance of a definite pre-existing
position, the CHSH inequality would not be violated.
It is violated.
The potentialities encoded in $\psi$ cannot therefore be epistemic in
the relevant sense---they are not ignorance about something that is
already actual.

Second, Margenau argued as early as 1954 \cite{margenau_1954_interpretations}
that quantum states describe \emph{latencies}---real dispositional
properties of physical systems that are not reducible to the statistics
of measurement outcomes.
Bunge \cite{bunge_1979_ontology} developed this into a systematic
ontology of propensities, arguing that quantum probabilities are
grounded in real physical tendencies of systems, not in subjective
degrees of belief or objective frequencies over ensembles.
On both accounts, the potentialities are \emph{ontic}: they are
features of the world that ground the Born-rule statistics, not
descriptions of our epistemic state with respect to an already-actual
world.


\footnote{Barandes \cite{barandes_2023_stochastic} shows that the
full empirical content of quantum mechanics can be recovered from
\emph{indivisible stochastic processes}---transition probabilities
satisfying a generalised Markov condition---without positing a
wavefunction or Hilbert space as primitive ontology. Potentialities
become transition probabilities between configurations, arguably
closer to the Aristotelian notion of \dunamis\ as process than the
static wavefunction picture. The Barandes framework is too recent to evaluate
fully, but illustrates that the ontological layer of quantum
mechanics remains genuinely open}

The potentiality reading we are defending is therefore not epistemic
Copenhagen in disguise.
It is a realist position that takes the wavefunction to represent
real features of the world---features that are genuinely potential
rather than actual, in the Aristotelian sense, prior to measurement.
The difference from standard collapse interpretations is that we do
not treat potentiality as a temporary state of ignorance to be
resolved by measurement into something fully actual and classical.
Potentiality is a permanent and irreducible mode of being for
quantum systems between interactions---it is what they \emph{are},
not what we \emph{don't know} about them.


% Como \footnote{} al final del párrafo que lista Copenhagen/GRW/Everett/Bohm
% Section .............................................................................................................
\section{Entanglement, Structural Realism, and Conclusion}
\label{section:structural}

% Subsection ..........................................................................................................
\subsection{Entanglement as Irreducibly Relational Potentiality}

The potentiality reading connects naturally to ontic structural realism (OSR) \citep{ladyman_2007_everything, french_2014_structure}. OSR proposes that what is real
are structures---patterns of relations---rather than objects with intrinsic natures, motivated in part by the permutation symmetry of
identical particles \citep{french_1988_indiscernibles} and by the non-locality revealed by Bell inequality violations \cite{bell_1964_epr, aspect_1982_bell, hensen_2015_loophole}.

\medskip

Consider the spin-singlet state of two electrons,
\begin{equation}
  \ket{\Psi^-} = \frac{1}{\sqrt{2}}\bigl(\ket{\uparrow}_A\ket{\downarrow}_B - \ket{\downarrow}_A\ket{\uparrow}_B\bigr).
  \label{eq:bell}
\end{equation}

This state cannot be written as a product $\ket{\phi}_A\otimes\ket{\chi}_B$. Neither particle $A$ nor $B$ has a definite spin potentiality on its own. On the Aristotelian reading, the potentialities are properties not of $A$ and $B$ separately but of the \emph{pair} $AB$ jointly---an irreducibly relational structure.

\medskip

Bell's theorem \citep{bell_1964_epr} shows that no local hidden-variable theory can reproduce the quantum correlations. The CHSH inequality
\begin{equation}
  \bigl|E(a,b) - E(a,b') + E(a',b) + E(a',b')\bigr| \leq 2
  \label{eq:chsh}
\end{equation}
is violated by quantum mechanics up to $2\sqrt{2}$ (the Tsirelson bound), confirmed experimentally \citep{aspect_1982_bell, hensen_2015_loophole}. Entanglement is not a correlation between independently real objects. It is a structure of irreducibly joint potentiality with no local decomposition---precisely the structural realist's claim that relations are ontologically prior to relata.

% Subsection ..........................................................................................................
\subsection{Potentiality Realism vs.\ Ontic Structural Realism}

The question arises whether the Aristotelian potentiality reading simply \emph{is} a form of OSR. We suggest it occupies a distinct position. OSR in its strongest form eliminates objects entirely in favour of structures \citep{french_1988_indiscernibles}. The Aristotelian reading preserves objects---the electron is a genuine entity---but denies that objects are fully actual with respect to all their properties prior to measurement.

\medskip

This is a form of \emph{potentiality realism}: realism about both objects and their potentialities, with actuality restricted to what has been, or is being, measured. Objects are real; their potentialities are real; their actualities are contingent on measurement interactions.

\medskip

\TODO{Engage with Esfeld and Deckert's minimalist ontology \citep{esfeld_2017_minimalist}: points in a structure with no intrinsic properties. Argue potentiality realism is distinct: it accepts intrinsic properties \emph{in potentia}, not \emph{in actu}.}

% Subsection ..........................................................................................................
\subsection{Conclusion}

The classical calculation of the hydrogen atom's radiative collapse is more than a historical curiosity. It is a precise, mathematically impeccable derivation---from Newton, Coulomb, and Larmor, all independently confirmed in their domains---that produces an empirically false result. And the reason it is false is not a flaw in the equations. It is a flaw in the ontology they presuppose.

\medskip

We have argued for three connected claims. First, every physical model carries an ontological layer analytically distinct from its predictive-mathematical content; diagnosing which layer has failed is a non-trivial philosophical task that physics alone cannot perform. Second, Bohr's 1913 resolution was an explicit ontological stipulation---he declared the classical ontology inapplicable without fully replacing it---and the residue of this incompleteness persists in the renormalisation programme of QED. Third, Aristotle's potentiality/actuality distinction, recovered by Heisenberg in his 1958 essay, provides a more adequate starting ontology for quantum mechanics than either classical particle realism or strict instrumentalism. Finally we discuss entanglement, which read through this lens, is a structure of irreducibly joint potentiality that connects naturally to ontic structural realism while preserving a genuine role for objects.

\medskip

The lesson is general. Scientific progress is not only the accumulation of more accurate predictions. It is the successive replacement of ontologies---each adequate to a domain, each carrying implicit metaphysical commitments, each eventually confronted with phenomena that force not just new equations but a new picture of what there is.

\medskip

Sixteen picoseconds. That is how long a classical hydrogen atom would survive. That it does not survive that way, and that we needed a century of physics and philosophy to understand why, says something important about the relationship between mathematics, experiment, and the picture of the world that any physical theory always---whether acknowledged or not---carries with it.

% To:
\section{Appendix: Bell Inequalities and the Tsirelson Bound}
\label{app:bell}

This note collects the mathematical background behind the claims made in Sections~4 and~5 regarding Bell's theorem and its consequences. It is intended as a self-contained reference for readers unfamiliar with the formalism.

%% -----------------------------------------------------------
\subsection*{1. The EPR argument and the problem of correlations}

Einstein, Podolsky, and Rosen argued in 1935 \cite{einstein_1935_epr} that if two spatially separated particles produce perfectly anti-correlated measurement outcomes without any possible causal connection between the measurements, then the outcome of each measurement must have been determined in advance by some property carried locally by each particle---a \textit{hidden variable}. Quantum mechanics, they argued, is incomplete because it does not specify this variable.

\medskip

Bell \cite{bell_1964_epr} converted this from a philosophical dispute into a precise mathematical question: what statistical correlations are compatible with \textit{any} local hidden-variable theory, regardless of its internal mechanism?

%% -----------------------------------------------------------
\subsection*{2. Bell's original inequality and Pearson correlations}

Suppose two particles, $A$ and $B$, are produced together and travel to separate detectors. Each detector can be oriented in one of several directions, and each measurement yields an outcome $\pm1$. Let $\lambda$ be a hidden variable (or set of variables) distributed according to some probability density $\rho(\lambda)$, unknown but fixed at the moment the particles are produced. The local hidden-variable assumption is that the outcome at $A$ depends only on the local setting $a$ and on $\lambda$---not on what $B$ does---and symmetrically for $B$:
\begin{equation}
  A = A(a,\lambda) \in \{-1,+1\}, \qquad
  B = B(b,\lambda) \in \{-1,+1\}.
  \label{eq:lhv}
\end{equation}

The correlation between the outcomes at settings $a$ and $b$ is then the \textit{Pearson correlation coefficient} of the two $\pm1$ random variables $A$ and $B$:
\begin{equation}
  E(a,b) = \int \rho(\lambda)\, A(a,\lambda)\, B(b,\lambda)\, \mathrm{d}\lambda.
  \label{eq:pearson}
\end{equation}

This is identical to the standard Pearson product-moment correlation $r_{AB} = \langle AB \rangle / (\sigma_A \sigma_B)$, which here reduces to $\langle AB \rangle$ because $\sigma_A = \sigma_B = 1$ for $\pm1$ variables. $E(a,b) = +1$ means perfect agreement, $E(a,b) = -1$ means perfect anti-correlation, and $E(a,b) = 0$ means statistical independence.

\medskip

Bell showed that for any three measurement directions $a$, $b$, $c$, the correlations defined by~(\ref{eq:pearson}) must satisfy:
\begin{equation}
  |E(a,b) - E(a,c)| \leq 1 + E(b,c).
  \label{eq:bell-supp}
\end{equation}

The proof is pure arithmetic: since $A(b,\lambda)B(b,\lambda) = \pm1$, one has $|A(a,\lambda)B(b,\lambda) - A(a,\lambda)B(c,\lambda)| \leq 1 - A(b,\lambda)B(b,\lambda)A(b,\lambda)B(c,\lambda)$ pointwise in $\lambda$, and integrating over $\rho(\lambda)$ yields (\ref{eq:bell-supp}). No assumption about quantum mechanics enters---the inequality holds for \textit{any} distribution $\rho(\lambda)$ and any $\pm1$-valued functions $A(a,\lambda)$, $B(b,\lambda)$.

%% -----------------------------------------------------------
\subsection*{3. The CHSH inequality}

Clauser, Horne, Shimony, and Holt (1969) \cite{clauser_1969_chsh} reformulated Bell's result in a form better suited to real experiments, which suffer from detector inefficiencies and finite statistics. The four authors' initials give the inequality its name: \textbf{CHSH}.

\medskip

The setup uses two settings per detector---$a$ and $a'$ for particle $A$, and $b$ and $b'$ for particle $B$---yielding four correlation functions. The CHSH quantity $S$ is defined as:

\begin{equation}
  S = E(a,b) - E(a,b') + E(a',b) + E(a',b').
  \label{eq:S}
\end{equation}

The letter $S$ stands for the \textit{sum} (or \textit{score}) of these four Pearson correlations, combined with a specific sign pattern chosen so that any local hidden-variable model is forced into an algebraic contradiction. The CHSH inequality states:
\begin{equation}
  |S| \leq 2.
  \label{eq:chsh-supp}
\end{equation}

The proof is again arithmetic: since each $A(\cdot,\lambda), A'(\cdot,\lambda), B(\cdot,\lambda), B'(\cdot,\lambda) \in \{-1,+1\}$, the combination $AB - AB' + A'B + A'B' = A(B-B') + A'(B+B')$ takes values in which one of $(B - B')$ and $(B + B')$ is always zero and the other is $\pm2$, so the whole expression is $\pm2$ pointwise in $\lambda$, and therefore $|S| \leq 2$ after integrating over $\rho(\lambda)$.

%% -----------------------------------------------------------
\subsection*{4. Quantum mechanics violates the CHSH bound}

For the spin-singlet state $|\Psi^-\rangle = \frac{1}{\sqrt{2}} (|\uparrow\rangle_A|\downarrow\rangle_B - |\downarrow\rangle_A|\uparrow\rangle_B)$, quantum mechanics predicts:
\begin{equation}
  E_{\mathrm{QM}}(a,b) = -\cos(\theta_{ab}),
  \label{eq:eqm}
\end{equation}

where $\theta_{ab}$ is the angle between the measurement directions $a$ and $b$ in the Bloch sphere. Note that (\ref{eq:eqm}) is still a Pearson correlation between $\pm1$ outcomes---the quantum calculation gives the same object as (\ref{eq:pearson}), it just predicts a \textit{different} value for it than any local hidden-variable model can.

\medskip

Choosing the optimal angles $a = 0\circ$, $a' = 90\circ$, $b = 45\circ$, $b' = 135\circ$ (so that each pair of settings subtends $45\circ$):
\begin{equation}
  E(a,b) = E(a',b) = E(a',b') = -\cos 45\circ = -\tfrac{1}{\sqrt{2}}, \qquad  E(a,b') = -\cos 135\circ = +\tfrac{1}{\sqrt{2}},
\end{equation}

giving:
\begin{equation}
  S_{\mathrm{QM}}
  = -\tfrac{1}{\sqrt{2}} - \tfrac{1}{\sqrt{2}} - \tfrac{1}{\sqrt{2}} - \tfrac{1}{\sqrt{2}}
  = -2\sqrt{2} \approx -2.828.
\end{equation}

The magnitude $|S_{\mathrm{QM}}| = 2\sqrt{2} > 2$ violates the CHSH bound~(\ref{eq:chsh-supp}). This violation was confirmed experimentally by Aspect, Grangier, and Roger in 1982 \cite{aspect_1982_bell}, and without loopholes by Hensen et al.\ in 2015 \cite{hensen_2015_loophole}. 

%% -----------------------------------------------------------
\subsection*{5. The Tsirelson bound: why not stronger?}

The CHSH bound for classical correlations is $|S| \leq 2$. Quantum mechanics reaches $|S| = 2\sqrt{2} \approx 2.83$. But why does quantum mechanics stop there? Logically, $|S|$ could be as large as $4$ (if all four Pearson correlations had magnitude 1 simultaneously). Is there a deeper reason why the quantum world is \textit{more} non-local than any classical model but \textit{less} non-local than it could logically be?

\medskip

Boris Tsirelson (also transliterated Cirel'son) answered this in 1980 \cite{tsirelson_1980_bound}: \textit{$2\sqrt{2}$ is the maximum value of $|S|$ achievable by any quantum state and any quantum measurements, in any Hilbert space of any dimension.} This is the \textbf{Tsirelson bound}.

\medskip

To see intuitively why it is $2\sqrt{2}$ and not $4$: the four Pearson correlations $E(a,b)$, $E(a,b')$, $E(a',b)$, $E(a',b')$ cannot all equal $\pm1$ simultaneously in quantum mechanics, because the measurements do not commute. Choosing $a$ and $b$ to maximise $E(a,b)$ forces the other correlations into a configuration that cannot simultaneously be maximised---the non-commutativity of quantum observables introduces a geometric constraint that limits the total score $S$.

\medskip

The algebraic proof uses the operator version of $S$:
\begin{equation}
  \hat{S} = \hat{A}\otimes(\hat{B}+\hat{B}') + \hat{A}'\otimes(\hat{B}-\hat{B}'),
\end{equation}

where $\hat{A}, \hat{A}', \hat{B}, \hat{B}'$ are observables with eigenvalues $\pm1$ (Hermitian operators with $\hat{A}^2 = \mathbb{1}$, etc.). Computing $\hat{S}^2$ and using $\hat{A}^2 = \hat{A}'^2 = \hat{B}^2 = \hat{B}'^2 = \mathbb{1}$:
\begin{equation}
  \hat{S}^2 = 4\,\mathbb{1} + [\hat{A},\hat{A}']\otimes[\hat{B},\hat{B}'].
  \label{eq:S2}
\end{equation}
The operator norm of a commutator $[\hat{X},\hat{Y}]$ for $\pm1$ observables satisfies $\|[\hat{X},\hat{Y}]\| \leq 2$, so $\|\hat{S}^2\| \leq 4 + 4 = 8$, giving $\|\hat{S}\| \leq 2\sqrt{2}$. The bound is tight: the singlet state with $45\circ$ settings achieves it exactly.

%% -----------------------------------------------------------
\subsection*{6. The intermediate bound and its philosophical meaning}

The existence of an intermediate bound---classical $\leq 2$, quantum $\leq 2\sqrt{2}$, logical maximum $= 4$---is not an accident, and its philosophical significance goes beyond confirming quantum mechanics.

\medskip

Popescu and Rohrlich noted in 1994 \cite{popescu_1994_pr} that one can define hypothetical theories, now called \textit{PR-boxes}, that respect the no-signalling principle (no faster-than-light communication) but saturate $|S| = 4$. Such theories are logically consistent and non-signalling, yet they would allow correlations far stronger than anything quantum mechanics predicts. The fact that nature chooses $2\sqrt{2}$, and not $4$, is therefore a \textit{non-trivial physical fact} that is not explained by no-signalling alone. Why this particular number?

\medskip

Several research programmes have tried to derive $2\sqrt{2}$ from information-theoretic principles---non-trivial communication complexity \cite{brassard_2006_limits}, information causality \cite{pawlowski_2009_information}, and local orthogonality \cite{fritz_2013_local}---but no single principle has been universally accepted as the explanation.

\medskip

For our purposes, this is directly relevant to the potentiality reading defended in Section~5. The Tsirelson bound shows that the joint potentialities encoded in the quantum state are not arbitrary: they obey a precise algebraic structure (the Hilbert space formalism, and specifically the constraint~(\ref{eq:S2}) on operator norms) that limits \textit{how strongly} the potentialities of two entangled systems can be correlated.

\medskip

The entanglement is not a free, unlimited non-local resource---it is a structured form of joint potentiality, constrained by the algebra of quantum observables. This is consistent with potentiality realism: the potentialities are \textit{real} (they exceed classical bounds), but they are \textit{structured} (they obey the Tsirelson bound rather than the logical maximum).

\medskip

The gap between $2\sqrt{2}$ and $4$ is, in this sense, a quantitative signature of the ontological structure of quantum potentiality---a structure that neither classical ontology nor pure logic would have predicted, and that remains only partially understood.


\appendix

\bibliographystyle{unsrt}
\bibliography{3_bibliography/references}

\end{document}